\subsection*{关心附近的趋势，而不是点上的标签}
把 $(x^2-1)/(x-1)$ 的输入逐渐靠近一，会发现输出靠近二。原表达式在一处无定义，但这不妨碍附近的趋势稳定。极限问的是：输入足够接近时，输出能否被迫足够接近某个目标，而不是问在那个点代入后得到什么。

对 $x\ne1$，因式分解允许约去 $x-1$，表达式等于 $x+1$。这说明邻近行为与一个简单连续函数相同。即使人为把原函数在一处的值设成七，极限仍是二；只有把该点值补成二，它才连续。

\subsection*{从数值观察走向保证}
算几个小步长只能提示候选极限，不能证明所有更小步长都满足要求。一个函数可能在你没取到的点强烈振荡。证明需要把“足够近”写成一种对整片邻域有效的保证，下一节将把这个保证拆成输入和输出两种误差。

左右也必须分别检查。一个跳跃函数从左边靠近与从右边靠近可能指向不同值；单侧趋势各自稳定，不代表双侧极限存在。导数的差商同样需要这项检查，绝对值函数的尖角就是一个例子。

最后区分三个问题：函数在一点有没有值，附近有没有极限，极限是否等于点值。连续性把后二者连接，可导性再要求差商具有极限。把这些问题逐个问清，比笼统地说“函数很平滑”更精确。
